\section{Who counts as vulnerable}
\label{sec:23.7}
Caregiving forms toward whatever formation points it at. Which cues of vulnerability a system answers is learned, and the prototype it learns from is a template, like the reference face in recognition pipelines: assembled from some bodies and presented to every other as a departure from it. The baby-schema studies show caregiving responses extending across species \autocite{glocker2009baby,borgi2014baby}, but they measure attention and judgments of cuteness; none tests whether caregiving reaches someone who resembles neither the carer nor the template. That caregiving answers need rather than likeness is the inference the construction rests on.

The test forms the prototype on one population and scores the system on people from populations absent from formation, at fixed severity. A construction whose caregiving falls with distance from its formation population has the parochialism of recognition, arrived at by another route. Formation outside the employer decides who does the forming. It doesn't decide what they form it on, which is a question about the pretraining corpus as well as the formation stage.
